\section{Анализ на изискванията и потребителските сценарии}

След въвеждането на общия контекст и теоретичната рамка е необходимо системата да бъде описана от гледна точка на конкретните изисквания, които тя трябва да удовлетворява. Тази стъпка е важна, защото осигурява връзката между предметната област и реалната имплементация. Изискванията не са само списък от желани функции; те са критерий, спрямо който могат да бъдат оценени архитектурата, структурата на кода и степента, в която приложението изпълнява предназначението си.

В случая на онлайн магазин за компютърна и офис техника изискванията се формират както от типичните процеси в електронната търговия, така и от конкретния начин, по който проектът е структуриран в сорс кода. Анализът показва, че системата е ориентирана към сравнително малка, но функционално завършена среда, в която регистрираният потребител може да управлява своята информация, да въвежда фирми и продукти и да преминава през основния цикъл на покупка.

\subsection{Основни участници в системата}

В най-общ вид могат да бъдат разграничени три типа участници, които влияят върху поведението на системата.

\subsubsection{Посетител без активна сесия}

Това е потребител, който достъпва приложението без да е логнат. В този режим системата следва да предоставя начална страница, страница за регистрация и страница за вход. Посетителят няма право да разглежда профилни данни, да добавя компании, да въвежда продукти или да използва количката. От гледна точка на реализирания код тази логика е постигната чрез проверки за наличие на атрибут \texttt{email} в сесията.

\subsubsection{Потвърден регистриран потребител}

Това е основният активен участник в системата. След потвърждение на регистрацията чрез електронна поща и успешен вход той получава достъп до каталога, профила, списъка с потребители, фирмите, формите за добавяне на фирма и продукт, количката и историята на поръчките. Важно наблюдение е, че в текущия проект липсва отделна роля от тип администратор; много от действията, които в производствена среда биха били ограничени, тук са достъпни за всеки удостоверен потребител. Това е приемливо за учебен сценарий, но трябва да бъде отчетено като особеност на изискванията към настоящата имплементация.

\subsubsection{Външни инфраструктурни компоненти}

Освен човешките потребители системата взаимодейства и с външни компоненти. На първо място това е MySQL базата данни, която съхранява постоянната информация. На второ място това е SMTP сървърът, използван за изпращане на електронно писмо при регистрация. На трето място може да се разглежда и браузърният JavaScript код, който използва REST крайна точка за извличане на максимална цена на продукт. Макар тези компоненти да не са „актьори“ в класически смисъл, те формират съществена част от оперативния контекст на приложението.

\subsection{Основни потребителски цели}

Анализът на предметната област позволява да бъдат формулирани няколко основни потребителски цели:

\begin{itemize}[leftmargin=*]
    \item създаване на акаунт и потвърждение на електронна поща;
    \item успешно удостоверяване и работа в персонална сесия;
    \item въвеждане и поддържане на фирмена информация, асоциирана с потребителя;
    \item въвеждане на продукти с категория, цена, количество и изображение;
    \item разглеждане на каталог и избор на подходящ артикул;
    \item добавяне на избран продукт в количка с желано количество;
    \item финализиране на поръчка с адрес за доставка;
    \item преглед на минали поръчки и история на потребителската активност.
\end{itemize}

Тези цели са практическият еквивалент на функционалните изисквания. Те показват как системата се възприема от потребителя не като набор от контролери и таблици, а като последователност от очаквани резултати.

\subsection{Функционални изисквания}

Функционалните изисквания определят какви операции системата трябва да поддържа. В настоящата разработка те могат да бъдат формулирани по следния начин.

\subsubsection{Регистрация и потвърждение на акаунт}

Системата трябва да позволява създаване на нов потребителски профил чрез форма за регистрация. При успешно подадени данни следва да бъде създаден запис в базата, паролата да бъде съхранена в хеширан вид, а на посочения имейл да бъде изпратена връзка за потвърждение. До потвърждаване на акаунта достъпът до системата не трябва да бъде пълен. В текущата реализация това е постигнато чрез булево поле за потвърден статус и логика в потребителската услуга.

\subsubsection{Вход и изход от системата}

След потвърждаване на акаунта потребителят трябва да може да влезе в системата с имейл и парола. При успешен вход в сесията се съхранява идентификаторът на потребителя под формата на електронна поща, което позволява контролерите да проверяват дали сесията е активна. Системата трябва да поддържа и коректен изход, при който сесията се инвалидира.

\subsubsection{Управление на потребителски профил}

Удостовереният потребител трябва да може да преглежда и редактира основни лични данни. Това включва промяна на име, фамилия, имейл и парола. При смяна на паролата системата следва да изисква въвеждане на старата парола и да проверява коректността ѝ, преди да запази новата стойност в криптиран вид.

\subsubsection{Управление на фирми}

Потребителят трябва да може да въвежда фирми, които впоследствие да се използват при добавяне на продукти. За всяка фирма се поддържат име, ЕИК/идентификатор, ДДС номер, адрес и признак за регистрация. В контекста на текущия проект фирмата е логически свързана с конкретен потребител и се използва като източник на собственост за продуктите.

\subsubsection{Управление на продукти}

Системата трябва да позволява добавяне на нов продукт с наименование, цена, количество, категория, фирма и URL към изображение. Категорията следва да бъде избирана от предварително дефиниран набор от стойности, което улеснява контрола на данните и последващото филтриране в каталога. След добавяне продуктът трябва да бъде достъпен за визуализация в каталожната част на приложението.

\subsubsection{Каталог и преглед на детайли}

Всеки удостоверен потребител трябва да може да разглежда каталога на наличните продукти. Каталогът трябва да поддържа филтриране по категории, а за всеки продукт да се визуализират име, категория, цена и изображение. При избор на конкретен артикул системата трябва да предоставя детайлен изглед, от който да бъде възможно добавяне към количката.

\subsubsection{Количка и финализиране на поръчка}

След избор на продукт потребителят трябва да може да добави желано количество в количката, ако то е налично. При този процес системата трябва да проверява дали заявеното количество не надвишава текущата наличност. Количката трябва да визуализира избраните артикули, единичната и сумарната цена, както и да поддържа оформяне на поръчка чрез адрес за доставка и калкулация на крайна стойност.

\subsubsection{История на поръчките}

След приключване на покупката системата трябва да съхрани исторически запис за направената поръчка, включително дата, адрес и обща стойност. Потребителят трябва да може да разглежда тези записи в отделен екран и при необходимост да изчиства натрупаната история.

\subsubsection{Локализация и допълнителен REST интерфейс}

Системата трябва да поддържа базова смяна на езика между английски и български за част от интерфейса чрез локализирани съобщения. Освен това трябва да предлага поне една примерна REST крайна точка, която връща агрегирана информация в машинно четим формат. В текущия проект това е информацията за максималната цена на продукт.

\begin{table}[H]
    \centering
    \caption{Обобщение на основните функционални изисквания}
    \label{tab:functional-requirements}
    \begin{tabularx}{\textwidth}{L{3.0cm}YY}
        \toprule
        Идентификатор & Изискване & Реализация в проекта \\
        \midrule
        FR-1 & Регистрация на потребител & Форма за регистрация, запис в базата, хеширане на парола \\
        FR-2 & Потвърждение на акаунт & Имейл с линк и поле \texttt{isConfirmed} \\
        FR-3 & Вход и изход & Проверка на парола и управление на \texttt{HttpSession} \\
        FR-4 & Редакция на профил & Форма за профил и проверка на стара парола \\
        FR-5 & Добавяне на фирма & Отделен контролер и binding модел за фирмени данни \\
        FR-6 & Добавяне на продукт & Категории, фирма, количество, цена и URL изображение \\
        FR-7 & Разглеждане на каталог & Категорийни маршрути и Thymeleaf визуализация \\
        FR-8 & Добавяне в количка & Детайлен изглед, количество и проверка за наличност \\
        FR-9 & Финализиране на поръчка & Адрес, обща цена и създаване на история \\
        FR-10 & История на поръчките & Екран с исторически записи и операция за изчистване \\
        FR-11 & Локализация & \texttt{messages.properties}, \texttt{messages\_bg.properties}, interceptor \\
        FR-12 & REST услуга & Крайна точка за максимална цена на продукт \\
        \bottomrule
    \end{tabularx}
\end{table}

Таблица \ref{tab:functional-requirements} показва, че функционалностите на системата следват логична последователност от регистрация и вход към управление на данни, покупка и проследяване на резултата от тази покупка. Това е именно цикълът, който определя проекта като завършено макар и компактно електронно търговско приложение.

\subsection{Нефункционални изисквания}

Наред с конкретните функции, системата следва да отговаря и на редица нефункционални изисквания. Те не описват какво прави приложението, а как следва да го прави.

\subsubsection{Поддържаемост}

Кодът трябва да бъде организиран така, че логиката на системата да бъде лесна за проследяване. Това предполага разделяне по пакети, отделяне на контролери, услуги и repository компоненти, както и използване на отделни модели за вход, вътрешен трансфер и визуализация. Поддържаемостта е особено важна в дипломна разработка, тъй като документът трябва да може ясно да обясни структурата на проекта.

\subsubsection{Четим и последователен потребителски интерфейс}

Системата трябва да предлага интерфейс, който е достатъчно интуитивен за базови потребителски операции. Макар дизайнът да не е водещ изследователски фокус, от приложението се очаква да поддържа последователна навигация, разбираеми форми и ясни маршрути към основните действия.

\subsubsection{Коректност на входните данни}

Всички основни форми следва да бъдат защитени чрез валидация, така че очевидно невалидни стойности да не достигат до persistence слоя. В разглеждания проект това се реализира чрез анотации и обработка на резултата от bind операцията.

\subsubsection{Устойчиво съхранение на информацията}

Приложението трябва да използва релационна база данни, за да гарантира, че потребители, фирми, продукти, количка и история се съхраняват устойчиво и са достъпни между отделни изпълнения на системата. Използването на MySQL и Hibernate осигурява именно този тип устойчивост.

\subsubsection{Елементарна сигурност}

Необходимо е паролите да бъдат съхранявани хеширано, а достъпът до определени страници да бъде ограничен до удостоверени потребители. Макар това да не реализира пълноценна корпоративна схема за сигурност, то е минимално необходимо условие, за да бъде системата функционално и методологично коректна.

\subsubsection{Разширяемост}

Архитектурата следва да позволява добавяне на нови категории, нови екрани и нови услуги без фундаментално пренаписване на вече изградената основа. Именно поради това изборът на многослоен модел и Spring компоненти има не само настоящ, но и стратегически смисъл.

\subsection{Ключови потребителски сценарии}

За по-голяма яснота функционалните изисквания могат да бъдат представени и като сценарии. Това е полезно, защото показва не просто кои функции съществуват, а как те се комбинират в завършен работен поток.

\subsubsection{Сценарий 1: Регистрация и първи вход}

Посетителят отваря началната страница и избира регистрация. Той попълва име, фамилия, имейл и парола. След успешно изпращане системата създава нов потребител, хешира паролата и изпраща линк за потвърждение. След активиране на акаунта потребителят се връща към формата за вход, въвежда данните си и системата създава активна сесия.

\subsubsection{Сценарий 2: Добавяне на фирма и продукт}

След успешен вход потребителят отваря екрана за добавяне на фирма и въвежда необходимите атрибути. След това преминава към формата за нов продукт, избира една от наличните категории, посочва фирмата, задава цена, количество и URL на изображение. След запис продуктът става част от каталога.

\subsubsection{Сценарий 3: Разглеждане на каталог и покупка}

Потребителят разглежда каталога, филтрира по категория и избира конкретен продукт. Отваря детайлния екран, въвежда желано количество и го добавя в количката. Ако количеството е налично, системата актуализира наличността и създава или обновява запис в количката. След това потребителят отваря количката, въвежда адрес за доставка и потвърждава поръчката. Системата създава исторически запис и изчиства активната количка.

\subsubsection{Сценарий 4: Преглед на история и профил}

Във всеки момент удостовереният потребител може да отвори профила си, да редактира лични данни и да провери асоциираните компании. Той може да отвори и историята на поръчките, където да види натрупаните записи и при необходимост да ги изчисти.

\subsection{Ограничения, произтичащи от текущата постановка}

Анализът на изискванията разкрива и някои характерни ограничения. Липсата на отделни роли означава, че системата не разграничава административни от потребителски действия. Локализацията е частична и не обхваща изцяло всички текстове и категории. Някои операции, които в production среда биха били защитени чрез по-строга политика, тук са реализирани с по-прост модел. Същевременно тези ограничения не обезсмислят проекта; напротив, те очертават ясна граница между учебна система и пълноценна комерсиална платформа.

\subsection{Изводи}

Изискванията към системата показват, че проектът е замислен като компактно, но завършено уеб приложение, покриващо най-важните процеси на един онлайн магазин. Формулираните функционални и нефункционални критерии създават здрава основа за анализа на архитектурата и кода в следващите глави. Особено важно е, че изискванията са достатъчно конкретни, за да позволят последваща верификация, и достатъчно реалистични, за да направят системата представителна за практическо приложение в областта на електронната търговия.
